\documentclass{article}
\usepackage[UTF8]{ctex}
\usepackage{amsmath,amssymb,amsfonts}
\usepackage{enumitem}
\usepackage{geometry}
\geometry{a4paper, margin=1in}
% === 错题本排版优化（不参与题目标识） ===
\setlength{\parindent}{0pt}
\setlength{\parskip}{0.35em plus 0.10em minus 0.08em}
\setlength{\emergencystretch}{3em}
\setlist{itemsep=0.25em, topsep=0.40em, parsep=0pt}
\allowdisplaybreaks[2]
% === 错题本排版优化结束 ===
\begin{document}
\section*{合肥工业大学2026年高等代数真题}
\section*{102．合肥工业大学2026年高等代数试卷}

\subsection{题目1}
一．（12 分）计算 \(n\) 阶行列式

\[
D_{n}=\left|\begin{array}{cccc}
a_{1}-x_{1} & a_{2} & \cdots & a_{n} \\
a_{1} & a_{2}-x_{2} & \cdots & a_{n} \\
\vdots & \vdots & & \vdots \\
a_{1} & a_{2} & \cdots & a_{n}-x_{n}
\end{array}\right|
\]

其中 \(x_{1} x_{2} \cdots x_{n} \neq 0\) ．

\subsection{题目2}
二．（15 分）已知线性方程组

\[
\left\{\begin{array}{l}
x_{1}+x_{2}+x_{3}+x_{4}=-1 \\
3 x_{1}+2 x_{2}+4 x_{3}-x_{4}=0 \\
5 x_{1}+3 x_{2}+7 x_{3}-3 x_{4}=1 \\
a x_{1}+x_{2}+5 x_{3}+b x_{4}=c
\end{array}\right.
\]

有 3 个线性无关的解。

\begin{itemize}
  \item[（1）] 证明该线性方程组系数矩阵的秩为 2 。
  \item[（2）] 求参数 \(a, b, c\) 的值以及该线性方程组的通解．
\end{itemize}

\subsection{题目3}
三．（14 分）设 \(A\) 是数域 \(P\) 上的 \(n\) 阶可逆方阵，\(\alpha\) 与 \(\beta\) 为 \(P\) 上的 \(n\) 维列向量。

\begin{itemize}
  \item[（1）] 试用分块矩阵理论证明：\(\left|A-\alpha \beta^{\mathrm{T}}\right|=|A|\left(1-\beta^{\mathrm{T}} A^{-1} \alpha\right)\) 。
  \item[（2）] 当 \(\beta^{\mathrm{T}} A^{-1} \alpha=k \neq 1\) 时，求 \(A-\alpha \beta^{\mathrm{T}}\) 的逆矩阵。
\end{itemize}

\subsection{题目4}
四．（12 分）设 \(f(x)=x^{n}+a_{n-1} x^{n-1}+\cdots+a_{1} x+a_{0}\) 是数域 \(P\) 上的 \(n\) 次首一多项式，\(n \geq 1\) ，且 \(f(x)\)有 \(n\) 个根 \(x_{1}, x_{2}, \cdots, x_{n}\) ，重根按重数计算，\(c \in P\) 不是 \(f(x)\) 的根．证明：

\[
\sum_{i=1}^{n} \frac{1}{x_{i}-c}=-\frac{f^{\prime}(c)}{f(c)}
\]

其中 \(f^{\prime}(x)\) 表示 \(f(x)\) 的导数，即 \(f^{\prime}(x)=n x^{n-1}+(n-1) a_{n-1} x^{n-2}+\cdots+a_{1}\) ．

\subsection{题目5}
五．（15 分）设 \(V\) 是数域 \(P\) 上所有 3 阶对称矩阵关于通常矩阵加法与矩阵数乘构成的线性空间，考察\(V\) 的子空间

\[
U=\{A \in V \mid \operatorname{tr}(A)=0\}, W=\{k E \mid k \in P\} .
\]

其中 \(E\) 为3阶单位矩阵， \(\operatorname{tr}(A)\) 表示 \(A\) 的迹，证明：\(V=U \oplus W\) 。

\subsection{题目6}
六．（14 分）设 \(A\) 是复数域 \(\mathbb{C}\) 上的 2 阶方阵，记 \(W_{A}=\left\{Y \in \mathbb{C}^{2 \times 2} \mid A Y=Y A\right\}\) 。

\begin{itemize}
  \item[（1）] 证明：\(W_{A}\) 是 \(\mathbb{C}^{2 \times 2}\) 的子空间。
  \item[（2）] 讨论 \(W_{A}\) 的维数所有可能的值．
\end{itemize}

\subsection{题目7}
七．（15 分）设 \(n\) 阶实对称矩阵 \(A=\left(\begin{array}{cccc}0 & 1 & \cdots & 1 \\ 1 & 0 & \cdots & 1 \\ \vdots & \vdots & & \vdots \\ 1 & 1 & \cdots & 0\end{array}\right), n \geq 2\) ．\\
（1）求 \(A\) 的所有特征值与特征向量．\\
（2）求 \(n\) 阶正交矩阵 \(Q\) ，使得 \(Q^{\mathrm{T}} A Q\) 为对角矩阵。

\subsection{题目8}
八．（12 分）设 \(V\) 是数域 \(P\) 上的 \(n\) 维线性空间， \(\mathcal{A}\) 是 \(V\) 上的线性变换，满足 \(\mathcal{A}^{2}=\mathcal{O}\) ．证明：存在 \(V\)的一组基，使得 \(\mathcal{A}\) 在这组基下的矩阵为 \(\left(\begin{array}{cc}O & O \\ E_{r} & O\end{array}\right)\) ，其中 \(r=\operatorname{dim} \mathcal{A}(V)\) ．

\subsection{题目9}
九．（14 分）设 \(V\) 是实数域 \(\mathbb{R}\) 上的 \(n\) 维线性空间， \(\mathcal{A}\) 是 \(V\) 上的线性变换．

\begin{itemize}
  \item[（1）] 证明：如果 \(n\) 是奇数，则 \(\mathcal{A}\) 有1维不变子空间。
  \item[（2）] 证明：如果 \(\mathcal{A}\) 没有1维不变子空间，那么 \(\mathcal{A}\) 必有2维不变子空间。
\end{itemize}

\subsection{题目10}
十．（15 分）设 \(A\) 是复数域上的 6 阶方阵，\(A\) 的最小多项式为 \(d(\lambda)=\left(\lambda^{2}-2 \lambda+2\right)^{2}(\lambda-1)\) ，且 \(\operatorname{tr}(A)=6\) ，这里 \(\operatorname{tr}(A)\) 表示 \(A\) 的迹．

\begin{itemize}
  \item[（1）] 求 \(A\) 的特征多项式以及 Jordan 标准型．
  \item[（2）] 求 \(A\) 的伴随矩阵 \(A^{*}\) 的 Jordan 标准形．
\end{itemize}

\subsection{题目11}
十一．（12 分）设 \(A\) 是 \(n\) 阶实矩阵，证明：\(A\) 是正定矩阵的充分必要条件是存在 \(n\) 阶正定矩阵 \(B\) ，使得 \(A=B^{2}\) ．
\end{document}
